\chapter{Mathematical and Numerical Preliminaries}\label{chap:ch02_preliminaries}

% --- local symbols (minimal; not in the shared preamble) -------------------
\providecommand{\Cou}{C}        % Courant number
\providecommand{\nlat}{N_\lat}  % number of latitude rows
\providecommand{\nlon}{N_\lon}  % number of longitude columns

This chapter assembles the mathematical scaffolding the rest of the book stands on: the
spherical geometry every component shares, the discrete grids on which the continuous fields
live, the metric factors that the curvature of the sphere forces into every gradient and
divergence, the distinction between \emph{prognostic} variables (which the model time-steps) and
\emph{diagnostic} ones (which it reconstructs each step), and the family of time-integration
schemes the code actually uses --- with their stability limits stated honestly. Wherever a
quantity appears in the source it is pinned to its location with a code reference. The notation
fixed here is used unchanged throughout Parts~II--V.

\section{Spherical coordinates and the metric}

\model{} lives on a sphere of radius $\Erad$ (\code{EARTH\_RADIUS} $=6.371\times10^{6}\,\mathrm{m}$)
rotating at rate $\Om$ (\code{OMEGA} $=7.2921\times10^{-5}\,\mathrm{rad\,s^{-1}}$). A point is
addressed by longitude $\lon\in[0,2\pi)$, latitude $\lat\in[-\pi/2,\pi/2]$, and a vertical
coordinate (geometric depth $z$ in the ocean, a terrain-following $\sigma$ coordinate in the
multi-level atmosphere). The horizontal velocity is $\uvec=(\uu,\vv)$, with $\uu$ eastward and
$\vv$ northward, and $\ww$ is the vertical component.

\codref{chronos\_esm/config.py}{physical constants $\Erad,\Om,\grav$ and grid sizes}

The essential complication of a sphere is that the zonal arc length shrinks toward the poles.
An eastward displacement $\dd\lon$ subtends a physical distance $\Erad\cos\lat\,\dd\lon$, whereas
a northward displacement $\dd\lat$ always subtends $\Erad\,\dd\lat$. The line element is
\begin{equation}
  \dd s^2 = \Erad^2\cos^2\!\lat\;\dd\lon^2 \;+\; \Erad^2\,\dd\lat^2 \;+\; \dd z^2 .
  \label{eq:line-element}
\end{equation}
Every horizontal differential operator inherits the $\cos\lat$ factor from
\eqref{eq:line-element}. For a scalar $f$ and horizontal vector $\uvec$,
\begin{align}
  \grad f &= \left(\frac{1}{\Erad\cos\lat}\pd{f}{\lon},\;\frac{1}{\Erad}\pd{f}{\lat}\right),
  \label{eq:sph-grad}\\
  \diver\uvec &= \frac{1}{\Erad\cos\lat}\left[\pd{\uu}{\lon}
     + \pd{(\vv\cos\lat)}{\lat}\right],
  \label{eq:sph-div}\\
  \lap f &= \frac{1}{\Erad^2\cos^2\!\lat}\pdd{f}{\lon}
     + \frac{1}{\Erad^2\cos\lat}\pd{}{\lat}\!\left(\cos\lat\,\pd{f}{\lat}\right).
  \label{eq:sph-lap}
\end{align}
The $\cos\lat$ in the denominators is the source of the polar singularity that haunts every
latitude--longitude model: as $\lat\to\pm\pi/2$, a fixed zonal grid increment maps to a vanishing
physical distance, so $\partial/\partial\lon$ and $\partial^2/\partial\lon^2$ amplify any
small-scale zonal structure without bound. Section~\ref{sec:metric-discrete} shows exactly how
the two grids of \model{} cope --- one by capping the metric, the other by working in a basis in
which the singularity never appears.

The planetary rotation enters through the Coriolis parameter
\begin{equation}
  \Cor = 2\Om\sin\lat, \qquad \beq \equiv \pd{\Cor}{y} = \frac{2\Om\cos\lat}{\Erad},
  \label{eq:coriolis}
\end{equation}
the second of which is the $\beta$ that controls Rossby propagation and the western
intensification of ocean gyres. In the ocean driver $\Cor$ is built directly from a uniform
latitude vector, \code{f\_cor = 2*OMEGA*sin(lat\_rad)}, and is floored away from zero near the
equator (to $|\Cor|\ge 10^{-5}\,\mathrm{s^{-1}}$) so that the geostrophic/Stommel division by
$\Cor$ stays finite there.

\codref{chronos\_esm/ocean/veros\_driver.py}{$\Cor$ from latitude, equatorial floor in \code{step\_ocean}}

\section{The model grids}

\model{} runs at spectral truncation \textbf{T31} by default, an option set by the environment
variable \code{CHRONOS\_RESOLUTION} (\code{T31} or \code{T63}). At T31 the horizontal grid is
$\nlon\times\nlat = 96\times48$, i.e.\ a nominal $3.75^\circ$ spacing.

\codref{chronos\_esm/config.py}{\code{RESOLUTION}, \code{NLAT=48}, \code{NLON=96}}

\subsection{The ocean grid: latitude--longitude, $z$-level}

The ocean is discretized on a regular latitude--longitude mesh with $96$ cells in longitude and
$48$ in latitude, the row centres laid out by \code{linspace(-90, 90, ny)}. Longitude spacing is
uniform, $\Delta\lon = 2\pi/\nlon$; the corresponding physical zonal spacing is latitude
dependent,
\begin{equation}
  \Delta x(\lat) = \Erad\cos\lat\,\Delta\lon, \qquad
  \Delta y = \frac{\pi\,\Erad}{\nlat}\;\;(\text{uniform}),
  \label{eq:ocean-dxdy}
\end{equation}
exactly the metric of \eqref{eq:line-element}. The driver receives $\Delta x$ as an array (it
varies down the column) and $\Delta y$ as a scalar.

Vertically the ocean carries $N_z=15$ levels on a \emph{stretched} grid, not a uniform one. A
uniform $5000/15\approx 333\,\mathrm{m}$ surface layer would give the sea surface an enormous
heat capacity and a sluggish response to fluxes, because surface forcing is applied as
$\text{flux}/(\dens_0 c_p\,\Delta z_0)$. Instead the interfaces are placed at
\begin{equation}
  z_i = \dep_{\mathrm{tot}}\left(\frac{i}{N_z}\right)^{\!s}, \quad i=0,\dots,N_z,\qquad
  \dep_{\mathrm{tot}}=5000\,\mathrm{m},\;\; s=1.7,
  \label{eq:ocean-vgrid}
\end{equation}
so that $\Delta z_k = z_{k+1}-z_k$ grows from a thin ($\approx 50\,\mathrm{m}$) surface layer to
a thick abyssal one, with layer-centre depths $\tfrac12(z_k+z_{k+1})$ used as the coordinate onto
which the World Ocean Atlas initial state is interpolated.

\codref{chronos\_esm/config.py}{\code{\_ocean\_vertical\_grid}, \code{OCEAN\_DZ}, \code{OCEAN\_DEPTH\_CENTERS}}

\begin{remark}
The ocean driver computes horizontal derivatives with plain centred differences of the form
$(\code{roll}(F,-1)-\code{roll}(F,1))/(2\Delta x)$ and treats $\Delta x,\Delta y$ as local
Cartesian spacings; the spherical \emph{convergence-of-meridians} term in the divergence,
$\partial(\vv\cos\lat)/\partial\lat$ of \eqref{eq:sph-div}, is approximated by the Cartesian
$\partial\vv/\partial\lat$ (no explicit $\cos\lat$ metric in the meridional divergence). This is a
known simplification; the diagnostic vertical velocity comment in the source flags it explicitly.
It is acceptable at this resolution but is one of the contributors to the spurious net meridional
transport that the per-latitude corrector of Chapter~\ref{chap:ch08_ocean_overturning} removes.
\end{remark}

\subsection{The atmosphere grids: spectral on a Gaussian mesh}

The atmosphere is \emph{spectral}: prognostic fields are represented not by their grid-point
values but by coefficients in a basis of spherical harmonics $Y_l^m(\lat,\lon)$, and the
``T31'' label is the triangular truncation total wavenumber $l\le 31$. Two atmospheres coexist
in the code.

\paragraph{Single-level (legacy).} The barotropic model in
\code{atmos/dynamics.py} is a hybrid pseudo-spectral scheme: it takes the FFT in longitude (exact
for the periodic direction) and uses fourth-order finite differences in latitude. Its grid is the
same regular $96\times48$ latitude--longitude mesh as the ocean, with the metric of
\eqref{eq:ocean-dxdy}.

\codref{chronos\_esm/atmos/spectral.py}{\code{compute\_gradients} --- FFT in $\lon$, 4th-order FD in $\lat$}

\paragraph{Multi-level (active).} The production atmosphere is Google's differentiable \dino{}
dycore (the engine behind NeuralGCM, \citealp{kochkov2024neuralgcm}), run at T31 with $24$
$\sigma$-levels and full spherical-harmonic transforms. Its horizontal grid is a \emph{Gaussian}
mesh: the latitude rows sit at the roots of the Legendre polynomial $P_{\nlat}(\sin\lat)$, chosen
so that Gauss--Legendre quadrature integrates the truncated harmonics exactly. Grid$\to$spectral
and spectral$\to$grid maps are \code{coords.horizontal.to\_modal} /
\code{to\_nodal}, and the eastward/northward winds are recovered from the modal state as
\begin{equation}
  \uu = \frac{(\widehat{\uu\cos\lat})_{\,0}}{\cos\lat}, \qquad
  \vv = \frac{(\widehat{\uu\cos\lat})_{\,1}}{\cos\lat},
  \label{eq:dino-uv}
\end{equation}
i.e.\ the dycore carries the momentum-like quantity $\uvec\cos\lat$ (regular at the poles) and
divides out $\cos\lat$ only for diagnostics --- the spectral route never divides by $\cos\lat$
internally, which is precisely why it has no polar instability.

\codref{chronos\_esm/atmos/dino\_atmos.py}{T31, $24$ $\sigma$-levels; \code{compute\_diagnostic\_state} for $\uu,\vv$}

\subsection{Metric terms in the discrete code}\label{sec:metric-discrete}

The single-level model must build the metric by hand. It forms the per-row zonal spacing
$\Delta x(\lat)=2\pi\Erad\cos\lat/\nlon$ and $\Delta y=\pi\Erad/\nlat$ directly. The polar
catastrophe implicit in the $1/\cos\lat$ factors of \eqref{eq:sph-grad}--\eqref{eq:sph-lap} is
dealt with by a hard floor on the cosine used in the \emph{dynamics}: rather than letting $\Delta x\to 0$ (it reaches $\approx 4\,\mathrm{m}$ in
the polar row when $\cos\lat$ is floored to $10^{-5}$), the metric is capped at its $80^\circ$
value,
\begin{equation}
  \cos\lat\big|_{\mathrm{dyn}} = \max\!\big(\cos\lat,\;\cos 80^\circ\big),
  \qquad \Delta x = \frac{2\pi\Erad\,\cos\lat\big|_{\mathrm{dyn}}}{\nlon},
  \label{eq:cos-floor}
\end{equation}
so the handful of pole-most rows integrate with the (still small) $80^\circ$ metric instead of
detonating. A separate, slightly different floor is applied to the diffusion $\Delta x$ so the
explicit hyperdiffusion stencil $\propto 1/\Delta x^4$ stays inside its stability bound. This is a
pragmatic fix, not a clean discretization: without it the mean sea-level pressure blew up at the
poles to $\sigma\sim 150\,\mathrm{hPa}$ against an observed $\sim 12\,\mathrm{hPa}$. The \dino{}
atmosphere needs no such floor because, as in \eqref{eq:dino-uv}, the singular $1/\cos\lat$ is
never formed.

\codref{chronos\_esm/atmos/dynamics.py}{\code{cos\_lat\_dyn}/\code{cos\_lat\_diff} floors in \code{step\_atmos}}

\section{Prognostic and diagnostic variables}\label{sec:prog-diag}

A central organizing idea is the split between \textbf{prognostic} variables --- those that carry
a time derivative and are advanced by the integrator --- and \textbf{diagnostic} variables ---
those reconstructed each step from the prognostic state by solving a (usually elliptic) balance.
The split is not cosmetic: it is what makes a step well-posed, and it determines what must be
saved in a checkpoint to restart bit-for-bit.

In the ocean, the \code{OceanState} carries eight fields, but only the three tracers are truly
prognostic in the time-stepping sense. Temperature $\Temp$, salinity $\Salt$, and dissolved
inorganic carbon are advanced by an explicit Runge--Kutta step (Section~\ref{sec:time}); density,
the velocities $(\uu,\vv,\ww)$, and the barotropic streamfunction $\psib$ are \emph{diagnosed}
each step from the tracers via the equation of state, thermal wind, continuity, and a Stommel
balance respectively. Table~\ref{tab:prog-diag} summarizes the split.

\codref{chronos\_esm/ocean/veros\_driver.py}{\code{OceanState}; tracers stepped, $\uvec,\ww,\dens,\psib$ diagnosed}

\begin{table}[ht]
\centering
\caption{Prognostic vs.\ diagnostic variables in the ocean and the active atmosphere.}
\label{tab:prog-diag}
\begin{tabular}{@{}lll@{}}
\toprule
Component & Prognostic (time-stepped) & Diagnostic (reconstructed) \\
\midrule
Ocean & $\Temp,\;\Salt,\;\text{DIC}$ & $\dens,\;\uu,\;\vv,\;\ww,\;\psib$ \\
Atmosphere (\dino) & $\vort,\;\delta,\;T',\;\ln p_s$, humidity & $\uu,\vv$ (from \eqref{eq:dino-uv}), $\Temp$, geopotential \\
\bottomrule
\end{tabular}
\end{table}

The diagnostic velocities follow from the linear equation of state
\begin{equation}
  \dens = \densz - \alpha(\Temp-\Temp_0) + \beta(\Salt-\Salt_0),
  \quad \densz=1025,\;\alpha=0.2,\;\beta=0.8,
  \label{eq:eos}
\end{equation}
followed by a thermal-wind/geostrophic relation for the baroclinic shear and a Stommel
streamfunction for the barotropic flow; these are derived in
Chapters~\ref{chap:ch04_ocean_equations} and~\ref{chap:ch06_ocean_numerics}. Because the
velocities are diagnostic, a restart need only store the tracers (plus $\psib$ as a warm start for
the elliptic solve); the rest is rebuilt on the first step. The same logic governs \dino{}
checkpoints, whose prognostic modal state is vorticity $\vort$, divergence $\delta$, temperature
variation $T'$, and log surface pressure $\ln p_s$; $\uu,\vv$ are diagnostic. (A subtle
consequence: restoring \code{sim\_time=None} on resume is mandatory, because the ImEx integrator
adds a tendency \emph{array} to the clock and a stored scalar would break that add.)

\codref{chronos\_esm/ocean/veros\_driver.py}{\code{equation\_of\_state} --- linear EOS, $\alpha,\beta$}

\section{Time-integration schemes}\label{sec:time}

\model{} uses different time integrators in different components, each chosen for the stiffness
and conservation demands of the equations it solves. We state each in the canonical
$\dt{\boldsymbol\Phi}=\mathcal{R}(\boldsymbol\Phi)$ form, where $\boldsymbol\Phi$ is the
prognostic state and $\mathcal{R}$ its tendency.

\subsection{Forward Euler --- the single-level atmosphere}

The legacy barotropic atmosphere advances with the simplest possible scheme,
\begin{equation}
  \boldsymbol\Phi^{n+1} = \boldsymbol\Phi^{n} + \Delta t\,\mathcal{R}(\boldsymbol\Phi^{n}),
  \label{eq:forward-euler}
\end{equation}
with $\Delta t_{\mathrm{atmos}} = 30\,\mathrm{s}$ (\code{DEFAULT\_DT\_ATMOS}). Forward Euler is
first-order accurate and only conditionally stable: for an oscillatory mode it is
\emph{unconditionally} amplifying in the strict linear sense, but at this tiny step the per-step
growth on the fastest physical waves is negligible. The relevant constraint is the
gravity-wave Courant number
\begin{equation}
  \Cou = \frac{c\,\Delta t}{\Delta x_{\min}}, \qquad c=\sqrt{\grav\dep},
  \label{eq:cfl}
\end{equation}
which at $\Delta t=30\,\mathrm{s}$ and the metric of \eqref{eq:cos-floor} is $\Cou\sim 0.06$ ---
far inside the stability region --- so the dominant error is dissipative drift, not blow-up.
Stability of the single-level model is bought elsewhere: by strong $\nabla^4$ hyperdiffusion that
damps the small scales (Chapter~\ref{chap:ch10_atmosphere_singlelevel}) and by velocity clamps
($|\uu|\le 80\,\mathrm{m\,s^{-1}}$) that act as a tripwire flagging instability rather than as
physics.

\codref{chronos\_esm/atmos/dynamics.py}{\code{step\_atmos} --- forward Euler, $\nabla^4$ filter, $80\,\mathrm{m\,s^{-1}}$ clamp}

\begin{remark}
A semi-implicit option exists but is \emph{not wired in}: \code{solve\_helmholtz} in
\code{atmos/spectral.py} solves $(\lap-\gamma)x=b$ via FFT-in-$\lon$ plus a tridiagonal solve in
$\lat$, with $\gamma=\Delta t^2 R T_0$ the implicit gravity-wave coupling. It would be needed only
if $\Delta t$ were raised enough to make $\Cou$ approach unity; at $30\,\mathrm{s}$ it buys
nothing, so the live path stays explicit. The book documents it (Chapter~\ref{chap:ch06_ocean_numerics})
because the same Helmholtz machinery underlies the elliptic solves elsewhere.
\end{remark}

\subsection{Fourth-order Runge--Kutta --- ocean tracers}

The ocean tracers are advanced with the classical four-stage Runge--Kutta scheme,
\begin{equation}
\begin{aligned}
  k_1 &= \mathcal{R}(\boldsymbol\Phi^n), &
  k_2 &= \mathcal{R}\!\big(\boldsymbol\Phi^n+\tfrac{\Delta t}{2}k_1\big), \\
  k_3 &= \mathcal{R}\!\big(\boldsymbol\Phi^n+\tfrac{\Delta t}{2}k_2\big), &
  k_4 &= \mathcal{R}\!\big(\boldsymbol\Phi^n+\Delta t\,k_3\big),
\end{aligned}
\qquad
  \boldsymbol\Phi^{n+1}=\boldsymbol\Phi^n+\frac{\Delta t}{6}\big(k_1+2k_2+2k_3+k_4\big),
  \label{eq:rk4}
\end{equation}
applied jointly to $(\Temp,\Salt,\text{DIC})$ with $\Delta t_{\mathrm{ocean}}=900\,\mathrm{s}$
(\code{DT\_OCEAN}, $15$ minutes). RK4 is fourth-order accurate and has a comfortably large stable
region; the binding constraint here is again advective CFL, $|\uvec|\Delta t/\Delta x\lesssim 1$,
easily satisfied by the slow ocean currents. The tendency $\mathcal{R}$ combines upwind horizontal
advection, upwind vertical advection (whose vertical Courant number $|\ww|\Delta t/\Delta z\sim
10^{-3}$ is tiny), Laplacian diffusion, and surface fluxes injected into the top wet cell.

\codref{chronos\_esm/ocean/veros\_driver.py}{\code{tendencies}, the four-stage RK4 update in \code{step\_ocean}}

After the RK4 update the tracers pass through three post-processing operators that are \emph{not}
part of the differential equation but are essential for a multi-century integration: an optional
scale-selective $\nabla^4$ Shapiro-type filter (off by default, removing only the $2\Delta x$
checkerboard mode); hard physical clamps $\Salt\in[30,38]\,\mathrm{psu}$, $\Temp\in[250,320]\,
\mathrm{K}$, and a sea-ice freezing floor $\Temp\ge 271.35\,\mathrm{K}$; and an additive
global-mean salinity renormalization to $34.7\,\mathrm{psu}$. The clamps are honest band-aids ---
the book quantifies where they bite (rarely, in well-tuned runs) in Chapter~\ref{chap:ch16_climatology}.

\subsection{Semi-implicit / ImEx --- the multi-level atmosphere}

The \dino{} dycore solves the dry primitive equations with an \emph{implicit--explicit}
Runge--Kutta scheme (specifically the strong-stability three-stage ImEx-RK of the dycore,
\code{imex\_rk\_sil3}). The state is split into a stiff linear part $L\boldsymbol\Phi$ --- the
fast gravity- and acoustic-like waves carried by the divergence/pressure/temperature coupling ---
and a slow nonlinear part $N(\boldsymbol\Phi)$:
\begin{equation}
  \dt{\boldsymbol\Phi} = \underbrace{L\,\boldsymbol\Phi}_{\text{implicit}}
  \;+\; \underbrace{N(\boldsymbol\Phi)}_{\text{explicit}} .
  \label{eq:imex-split}
\end{equation}
Each stage treats $L$ implicitly (solved exactly in spectral space, where $L$ is diagonal in the
vertical-mode/horizontal-wavenumber basis) and $N$ explicitly. The point of the split is that the
CFL limit \eqref{eq:cfl} no longer involves the fast gravity-wave speed $c$ (handled implicitly,
hence unconditionally stable) but only the much slower advective speed, so a far larger step is
admissible: $\Delta t = 20\,\mathrm{min}$ at T31. After each macro-step a horizontal hyperdiffusion
filter is applied as a separate operator,
\begin{equation}
  \widehat{\boldsymbol\Phi}_l \leftarrow \widehat{\boldsymbol\Phi}_l \,
  \exp\!\left[-\frac{\Delta t}{\tau}\left(\frac{l(l+1)}{L(L+1)}\right)^{\!p}\right],
  \label{eq:dino-diffusion}
\end{equation}
with diffusion order $p$ and time-scale $\tau=6\,\mathrm{h}$ chosen \emph{weak enough} that the
T31 baroclinic eddies ($l\sim10$--$15$) are allowed to grow --- a deliberate departure from the
single-level model, where eddies were spurious and diffusion was kept strong (see the
decoupling principle discussed in Chapters~\ref{chap:ch09_atmosphere_multilevel}
and~\ref{chap:ch18_jets_stormtracks}).

\codref{chronos\_esm/atmos/dino\_atmos.py}{\code{imex\_rk\_sil3}, \code{step\_with\_filters}, \code{horizontal\_diffusion\_step\_filter}}

\begin{table}[ht]
\centering
\caption{Time integrators used across \model{} (default T31).}
\label{tab:integrators}
\begin{tabular}{@{}llll@{}}
\toprule
Component & Scheme & Order & $\Delta t$ \\
\midrule
Single-level atmosphere & Forward Euler + $\nabla^4$ filter & 1 & $30\,\mathrm{s}$ \\
Ocean tracers           & Classical RK4                     & 4 & $900\,\mathrm{s}$ \\
Multi-level atmosphere  & ImEx-RK (\dino{}) + spectral filter & 3 & $20\,\mathrm{min}$ \\
Coupler exchange        & ---                               & --- & $3600\,\mathrm{s}$ \\
\bottomrule
\end{tabular}
\end{table}

The components are synchronized by a coupler that exchanges fluxes every
$\Delta t_{\mathrm{cpl}}=3600\,\mathrm{s}$ (\code{COUPLING\_INTERVAL}); within a coupling window
each component sub-cycles at its own step (Table~\ref{tab:integrators}). The coupling itself is a
loose, first-order-in-time exchange of time-mean surface fields, documented in
Chapter~\ref{chap:ch13_coupler_fluxes}.

\section{Notation conventions}

We close with the conventions used throughout. Operators $\grad,\diver,\curl,\lap$ are the
spherical forms \eqref{eq:sph-grad}--\eqref{eq:sph-lap} unless a Cartesian tangent-plane form is
explicitly noted. The material derivative is
$\Dt{f}=\pd{f}{t}+\uvec\!\cdot\!\grad f + \ww\,\pd{f}{z}$. Subscripts $0$ denote reference
constants ($\densz,\Temp_0,\Salt_0$); primes denote departures from a reference or zonal mean
($T'$). A circumflex $\widehat{(\cdot)}$ denotes a spectral (modal) coefficient. A superscript $n$
is the time level, so $\boldsymbol\Phi^{n+1}$ is the updated state. Transport is reported in
Sverdrups, $1\Sv=10^{6}\,\mathrm{m^3\,s^{-1}}$. Units are written in plain roman math, e.g.\
$\mathrm{m\,s^{-1}}$. Discrete grid operations name the array axes as the code does: for the ocean
state of shape $(N_z,\nlat,\nlon)$, \code{axis=2} is longitude, \code{axis=1} is latitude, and
\code{axis=0} is the vertical. Code references appear inline as
$\hookrightarrow$~\textit{Code:} pointers. Where a smooth surrogate replaces a non-differentiable
operation (e.g.\ $\softplus$ for a positive part, or a soft clip), it is named at the point of use;
the rationale is the subject of Chapter~\ref{chap:ch03_differentiable}.

\begin{example}
At T31 the polar-row physical zonal spacing is $\Delta x = \Erad\cos\lat\,(2\pi/\nlon)$. For the
row centred near $\lat=86.25^\circ$, $\cos\lat\approx 0.0654$, giving $\Delta x\approx 27\,
\mathrm{km}$ --- already $15\times$ smaller than the equatorial $\approx 417\,\mathrm{km}$. Push to
$\lat\to90^\circ$ and $\Delta x\to0$: this is the singularity \eqref{eq:cos-floor} caps, and the
reason the \dino{} dycore works in the harmonic basis instead.
\end{example}

\begin{exercise}
Using \code{config.py}, reconstruct the ocean vertical grid \eqref{eq:ocean-vgrid} for $N_z=15$,
$\dep_{\mathrm{tot}}=5000\,\mathrm{m}$, $s=1.7$. Tabulate the $15$ layer thicknesses $\Delta z_k$.
Confirm the surface layer is $\approx 50\,\mathrm{m}$ and the deepest $\approx 900\,\mathrm{m}$.
Then recompute with $s=1.0$ (uniform) and explain, using
$\partial\Temp/\partial t\sim\text{flux}/(\dens_0 c_p\Delta z_0)$, why the stretched grid gives a
more responsive sea surface temperature.
\end{exercise}

\begin{exercise}
Evaluate the gravity-wave Courant number \eqref{eq:cfl} for the single-level atmosphere using
$\dep=10\,\mathrm{km}$, $\grav=9.81\,\mathrm{m\,s^{-2}}$, $\Delta t=30\,\mathrm{s}$, and the
metric-floored polar $\Delta x$ from \eqref{eq:cos-floor} at $\lat=80^\circ$. How large could
$\Delta t$ become before $\Cou\to1$? Explain why, despite the headroom, the production atmosphere
still moved to the ImEx \dino{} dycore rather than simply raising $\Delta t$ here.
\end{exercise}

\begin{exercise}
Classify each field of \code{OceanState} (Table~\ref{tab:prog-diag}) as prognostic or diagnostic
by inspecting \code{step\_ocean}. For a checkpoint that must restart the run, which fields are
strictly necessary and which can be regenerated on the first step? Verify your answer against the
checkpoint-restart discussion in Chapter~\ref{chap:ch06_ocean_numerics}.
\end{exercise}

\begin{exercise}
The ImEx split \eqref{eq:imex-split} treats the fast linear waves implicitly. Write the linear
shallow-water gravity-wave system $\partial_t u=-g\partial_x h$, $\partial_t h=-H\partial_x u$ as
$\dt{\boldsymbol\Phi}=L\boldsymbol\Phi$, find the dispersion relation, and show that a fully
implicit (backward-Euler) treatment of $L$ is unconditionally stable while forward Euler
\eqref{eq:forward-euler} is not. Relate this to why \dino{} can take a $20\,\mathrm{min}$ step
where the explicit single-level model is pinned at $30\,\mathrm{s}$.
\end{exercise}
